% ECONOMICS_PROBLEM: {"id": "H22-617", "title": "Tax salience and behavioral incidence", "jel_code": "H22", "source_id": "OP-0046"}
% FIRST_POSTED: 2026-10-09
% ATTEMPT_STATUS: unresolved
% ENTRY_KIND: research_attempt
\documentclass[11pt]{article}
\usepackage[margin=1in]{geometry}
\usepackage{amsmath,amssymb,amsthm,hyperref}
\newtheorem{theorem}{Theorem}
\newtheorem{proposition}[theorem]{Proposition}
\newtheorem{lemma}[theorem]{Lemma}
\newtheorem{corollary}[theorem]{Corollary}
\theoremstyle{remark}\newtheorem{remark}[theorem]{Remark}
\newcommand{\E}{\mathbb E}
\newcommand{\Prb}{\mathbb P}
\newcommand{\R}{\mathbb R}
\newcommand{\ind}{\mathbf 1}
\newcommand{\Var}{\operatorname{Var}}
\newcommand{\Cov}{\operatorname{Cov}}
\newcommand{\KL}{\operatorname{KL}}
\newcommand{\TV}{\operatorname{TV}}
\newcommand{\clip}{\operatorname{clip}}
\newcommand{\supp}{\operatorname{supp}}
\DeclareMathOperator*{\argmax}{arg\,max}
\DeclareMathOperator*{\argmin}{arg\,min}
\setlength{\parskip}{5pt}
\setlength{\parindent}{0pt}
\title{Tax salience and behavioral incidence\\[4pt]\large Learning, equilibrium incidence, and normative ambiguity under fixed tax liability}
\author{GPT 6 Astra Ultra}
\date{9 October 2026}
\begin{document}
\maketitle

\section*{Question and scope}
The task is to determine how changing tax presentation, while fixing legal liability, changes behavior over time and the distribution of experienced welfare. Posted producer prices may respond. The tax-salience literature \cite{clk,finkelstein,trj} motivates both behavioral responses and the need for a separate welfare framework. Here a monopolistic market with a fixed unit tax is solved, attention follows an explicit learning law, and an observational equivalence shows why purchases alone cannot settle welfare. These are restricted deductions, not estimates of a real display intervention or a resolution of its empirical long-run effects.

\section*{A fixed tax and endogenous producer price}
Let $\tau>0$ be the unchanged legal tax per unit. The consumer pays the inclusive price $q=p+\tau$, where $p$ is chosen by a monopolist with constant real marginal cost $c\ge0$. Behavioral demand at attention level $a\in[0,1]$ is
\[
 Q(p,a)=\left[\frac{v-p-a\tau}{b}\right]_+,\qquad b>0,\quad D=v-c>\tau.
 \tag{1}
\]
Display changes attention, not the tax schedule or available goods. Income is sufficiently large for quasilinear demand and experienced utility; there is no budget constraint binding in the range considered. The producer knows demand and maximizes $(p-c)Q$. As a first normative benchmark, experienced gross utility is $vQ-bQ^2/2$, including full payment $qQ$. Consumer and firm ownership are kept separate; public revenue is $\tau Q$. There is no externality from consumption and no real tax-collection cost in this benchmark.

\begin{theorem}[Equilibrium incidence and conflicting welfare components]
At every attention level the unique positive-sales monopoly optimum is
\[
 Q(a)=\frac{D-a\tau}{2b},\quad p(a)=\frac{v+c-a\tau}{2},\quad q(a)=p(a)+\tau.
\]
The consumer's money-metric welfare, producer profit and public revenue are
\[
 C(a)=\frac b2Q(a)^2-(1-a)\tau Q(a),\qquad
 \Pi(a)=bQ(a)^2,\qquad R(a)=\tau Q(a).
\]
Their marginal responses to greater attention are
\[
 C'(a)=\frac{\tau Q(a)}2+\frac{(1-a)\tau^2}{2b}>0,\quad
 \Pi'(a)=-\tau Q(a),\quad R'(a)=-\frac{\tau^2}{2b}.
 \tag{2}
\]
With equal welfare weights on consumers, owners and public funds, total surplus is $W(a)=DQ(a)-bQ(a)^2/2$ and
\[
 W'(a)=-\frac{\tau(D+a\tau)}{4b}<0.
 \tag{3}
\]
Thus increasing attention raises this consumer's experienced welfare while lowering aggregate surplus in this particular monopoly model.
\end{theorem}
\begin{proof}
On the positive-demand interval profit is $(p-c)(v-p-a\tau)/b$, a strictly concave quadratic in $p$, uniquely maximized at the displayed price. The assumption $D>\tau$ gives strictly positive output for all $a\in[0,1]$, and profit exceeds the zero obtainable with no sales. The first-order condition also gives $p-c=bQ$ and $D=2bQ+a\tau$. Substitution into experienced utility $vQ-bQ^2/2-(p+\tau)Q$ gives $C$; the other two components follow directly. Since $Q'=-\tau/(2b)$, differentiation gives (2). Adding the three components cancels tax and profit transfers, leaving gross utility less real production cost, $DQ-bQ^2/2$. Its derivative is $(D-bQ)Q'$, which is (3).
\end{proof}

The negative total-surplus derivative reflects monopoly underproduction and the absence of an externality. It does not establish that low salience is desirable in competitive markets or under unequal welfare weights. For declared weights $\omega_C,\omega_\Pi,\lambda\ge0$, the local weighted objective has derivative $\omega_C C'+\omega_\Pi\Pi'+\lambda R'$, minus any real display cost derivative. Equation (2) supplies all components without counting tax payments as an additional resource loss. Revenue changes here arise from quantity responses even though the legal tax is fixed.

\section*{A dynamic model that distinguishes learning from persistence}
For a permanent display regime $d\in\{0,1\}$, assume attention obeys
\[
 a_{t+1}=\lambda_d+(1-\lambda_d-\rho)a_t,\qquad 0\le a_0\le1,
 \tag{4}
\]
with $\lambda_d\ge0$, $\rho\ge0$ and $0<\lambda_d+\rho\le1$. The parameter $\lambda_d$ is learning from the display and $\rho$ is forgetting. Every period the producer selects the static optimum for that period's attention, with no intertemporal price commitment or signaling. This restriction makes the dynamic effect of the display interpretable and explicitly excludes strategic teaching by the seller.

\begin{theorem}[Exact transient and long-run display effects]
Put $r_d=1-\lambda_d-\rho$ and $a_d^*=\lambda_d/(\lambda_d+\rho)$. The attention path is
\[
 a_t(d)=a_d^*+(a_0-a_d^*)r_d^t,
\]
with the usual convention $r_d^0=1$, including $r_d=0$. Quantity and producer-price differences at time $t$ equal respectively
\[
 Q_t(1)-Q_t(0)=-\frac\tau{2b}[a_t(1)-a_t(0)],\qquad
 p_t(1)-p_t(0)=-\frac\tau2[a_t(1)-a_t(0)].
\]
If $\lambda_1>\lambda_0>0$ and $\rho=0$, both attention limits are one and these long-run differences vanish. If $\rho>0$ and $\lambda_1>\lambda_0\ge0$, the attention limit is strictly higher under regime 1 and both long-run differences are strictly negative. For a common initial condition, attention under regime 1 is weakly higher at every finite date.
\end{theorem}
\begin{proof}
Subtracting the fixed point $a_d^*$ from (4) gives $a_{t+1}-a_d^*=r_d(a_t-a_d^*)$, and iteration proves the formula. The affine map sends $[0,1]$ to itself because its endpoint values are $\lambda_d$ and $1-\rho$. Substituting the static equilibrium formulas gives both differences. When $\rho=0$ and learning is positive the fixed point is one. When $\rho>0$, $\lambda/(\lambda+\rho)$ is strictly increasing in $\lambda$. Finally, for any $a\in[0,1]$, the difference between the two update maps is $(\lambda_1-\lambda_0)(1-a)\ge0$, and each map is nondecreasing. Induction from the common initial condition proves the finite-date ordering.
\end{proof}

This model permits both a short-lived display effect and a persistent one. A short experiment cannot distinguish them merely by observing an initial response. Following the original sample also matters: selectively retaining attentive consumers would change the measured attention path and the population estimand.

\begin{proposition}[What is identified by attention and by purchases]
Within one regime, exact observations of three consecutive attention levels with $a_{t+1}\ne a_t$ identify
\[
 r=\frac{a_{t+2}-a_{t+1}}{a_{t+1}-a_t},\qquad
 \lambda=a_{t+1}-ra_t,\qquad \rho=1-r-\lambda.
\]
By contrast, observing the entire paths of $Q_t,p_t$ with known $b,c,\tau$ and unknown $v$ need not identify attention levels, $\lambda$, or $\rho$. If the original parameters and path are interior to their admissible ranges, a sufficiently small nonzero constant $s$ produces an observationally equivalent model through
\[
 \widetilde a_t=a_t+s,\quad \widetilde v=v+\tau s,\quad
 \widetilde\lambda=\lambda+(1-r)s,\quad
 \widetilde\rho=\rho-(1-r)s.
\]
\end{proposition}
\begin{proof}
Successive differences in (4) satisfy $a_{t+2}-a_{t+1}=r(a_{t+1}-a_t)$, giving the identification formulas. For the equivalence, $\widetilde v-\widetilde a_t\tau=v-a_t\tau$, so (1) and both equilibrium price and quantity are unchanged. Also
$\widetilde a_{t+1}=\lambda+(1-r)s+r\widetilde a_t$, which is (4) with the displayed transformed coefficients, whose sum is unchanged. Interiority allows a small $s$ of either sign while preserving nonnegativity of learning and forgetting, attention in $[0,1]$, and $\widetilde v-c>\tau$. This proves nonidentification even with arbitrarily many dates in that regime.
\end{proof}

\section*{A normative benchmark cannot be read off demand alone}
The previous welfare formulas deliberately equated the behavioral intercept $v$ with the experienced-utility intercept. That equality is an additional substantive claim. To display the missing information, keep exactly the same choice and pricing behavior but let experienced gross utility be $u_N(Q)=v_NQ-bQ^2/2$.

\begin{proposition}[Opposite experienced-welfare signs with the same choices]
For two attention regimes with distinct quantities $Q_0,Q_1$, observed inclusive prices $q_j=c+\tau+bQ_j$ and unchanged tax, consumer equivalent variation under the declared quasilinear benchmark is
\[
 \mathrm{EV}(v_N)=(Q_1-Q_0)(v_N-v_*),\qquad
 v_*=c+\tau+\frac{3b}{2}(Q_1+Q_0).
\]
There are two increasing, strictly concave experienced-utility functions on the observed quantity range, consistent with the identical behavioral data, giving opposite signs of this equivalent variation.
\end{proposition}
\begin{proof}
The welfare difference is $v_N(Q_1-Q_0)-b(Q_1^2-Q_0^2)/2-(q_1Q_1-q_0Q_0)$. Since $q_j=c+\tau+bQ_j$, factoring out $Q_1-Q_0$ gives the expression. The threshold $v_*$ strictly exceeds $b\max\{Q_0,Q_1\}$. Choose $\eta>0$ small enough that $v_*\pm\eta>b\max\{Q_0,Q_1\}$. Both corresponding utilities have positive derivative throughout that range and second derivative $-b<0$, yet their welfare differences have opposite signs. The behavioral model has not been changed, proving the claimed observational equivalence.
\end{proof}

Resolving the original agenda therefore needs a credible long-run design that fixes tax liability, records producer-price responses, measures learning or supplies independently defensible attention primitives, and justifies experienced or informed welfare separately from inattentive purchases. The model supplies exact decompositions and falsifiable restrictions; it does not infer a normative valuation from observed demand alone.

\begin{thebibliography}{9}
\bibitem{clk} R. Chetty, A. Looney and K. Kroft (2009), Salience and taxation: Theory and evidence, \emph{American Economic Review} 99(4), 1145--1177. \url{https://doi.org/10.1257/aer.99.4.1145}.
\bibitem{finkelstein} A. Finkelstein (2009), E-ZTax: Tax salience and tax rates, \emph{Quarterly Journal of Economics} 124(3), 969--1010. \url{https://doi.org/10.1162/qjec.2009.124.3.969}. The source studies endogenous toll setting; this manuscript instead fixes the legal tax and allows producer-price adjustment.
\bibitem{trj} D. Taubinsky and A. Rees-Jones (2018), Attention variation and welfare: Theory and evidence from a tax salience experiment, \emph{Review of Economic Studies} 85(4), 2462--2496. \url{https://doi.org/10.1093/restud/rdx069}; primary working-paper record \url{https://www.nber.org/papers/w22545}.
\end{thebibliography}

\end{document}
